\subsection{Marco Te\'orico}

El dise\~no, desarrollo y evaluaci\'on de la plataforma web 3D se fundamenta en un marco te\'orico multidisciplinario que integra conocimientos de ingenier\'ia, dise\~no, psicolog\'ia cognitiva e inform\'atica gr\'afica. El proyecto no se limita a mostrar un modelo tridimensional: busca demostrar que una visualizaci\'on t\'ecnica interactiva puede mejorar la comprensi\'on espacial del hardware, sostener un rendimiento web aceptable y articular una experiencia de usuario m\'as clara que la documentaci\'on 2D est\'atica.

\paragraph{Hip\'otesis de usabilidad y carga cognitiva.} La visualizaci\'on 3D interactiva gestiona la carga cognitiva intr\'inseca mediante segmentaci\'on del contenido y reduce la carga extr\'inseca al externalizar operaciones de rotaci\'on mental que en soportes 2D deben resolverse internamente por el usuario. Esta hip\'otesis se validar\'a mediante tareas representativas, KPIs t\'ecnicos, SUS, NASA-TLX y observaci\'on \textit{Think-Aloud}.

El uso de una estrategia tipo MVP justifica la exclusi\'on intencional de caracter\'isticas no esenciales para la validaci\'on acad\'emica inmediata, como persistencia avanzada, instrumentaci\'on anal\'itica compleja o variantes de producto no integradas en la versi\'on final del prototipo.

\subsubsection{Teor\'ia de la Carga Cognitiva (Sweller, 1988; Sweller et al., 2019)}

La Teor\'ia de la Carga Cognitiva plantea que la memoria de trabajo tiene capacidad limitada y que el dise\~no de materiales e interfaces debe regular la demanda que impone sobre ella. Seg\'un Cowan (2001), esta capacidad operativa ronda unos pocos elementos simult\'aneos, por lo que la forma de presentar informaci\'on t\'ecnica incide directamente en la comprensi\'on y en la fatiga del usuario.

Sweller, van Merri\"enboer y Paas (2019) distinguen tres tipos de carga:

\paragraph{Carga intr\'inseca.} Corresponde a la complejidad inherente del contenido. En este proyecto es elevada porque el dron integra componentes estructurales, electr\'onicos y mec\'anicos con relaciones espaciales y funcionales entre s\'i.

\paragraph{Carga extr\'inseca.} Es la carga innecesaria generada por una presentaci\'on deficiente. En documentaci\'on 2D est\'atica, el usuario debe imaginar relaciones tridimensionales no visibles, lo cual incrementa la necesidad de rotaci\'on mental y consume recursos cognitivos que no est\'an dirigidos al objetivo t\'ecnico principal.

\paragraph{Carga germana.} Es el esfuerzo que s\'i contribuye a construir esquemas mentales \'utiles. En una interfaz 3D bien dise\~nada, los recursos cognitivos liberados por la reducci\'on de carga extr\'inseca pueden redirigirse a comprender mejor la funci\'on, ubicaci\'on y relaci\'on entre componentes.

En el contexto del prototipo, la interfaz 3D orbital y la consulta contextual de piezas buscan precisamente reducir la carga extr\'inseca y favorecer una carga germana orientada a comprensi\'on estructural.

\subsubsection{Navegaci\'on espacial e interacci\'on 3D}

Bowman et al. (2004) definen la interacci\'on 3D como aquella en la cual las tareas del usuario se ejecutan directamente en un contexto espacial tridimensional. Esto es relevante porque el valor del prototipo depende de que el usuario pueda explorar el ensamblaje sin perder contexto global mientras inspecciona detalles locales.

Los controles orbitales siguen principios de navegaci\'on espacial orientados a prevenir desorientaci\'on:
\begin{itemize}
    \item rotaci\'on alrededor de un foco de inter\'es estable;
    \item \textit{pan} y \textit{zoom} como operaciones separadas de la rotaci\'on;
    \item recentrado del inter\'es cuando se selecciona una pieza o se invoca una herramienta de inspecci\'on;
    \item continuidad de movimiento suficiente para evitar saltos bruscos de interpretaci\'on espacial.
\end{itemize}

\paragraph{Cognici\'on distribuida (Hutchins, 1995).} La interfaz no solo presenta informaci\'on: tambi\'en funciona como artefacto cognitivo. La c\'amara orbital externaliza parte de la rotaci\'on mental, los \textit{hotspots} y fichas de pieza operan como memoria externa y los modos anal\'iticos convierten procesos de inspecci\'on complejos en operaciones visuales guiadas.

\subsubsection{Heur\'isticas de usabilidad (Nielsen, 1994)}

Las heur\'isticas de Nielsen orientan la construcci\'on de la interfaz:
\begin{enumerate}
    \item \textbf{Visibilidad del estado del sistema.} La aplicaci\'on debe comunicar con claridad qu\'e pieza est\'a seleccionada, qu\'e modo visual est\'a activo y cu\'ando una herramienta de an\'alisis modifica la escena.
    \item \textbf{Correspondencia con el mundo real.} Los nombres de componentes, la organizaci\'on por categor\'ias y la l\'ogica de inspecci\'on deben alinearse con la terminolog\'ia t\'ecnica del dron.
    \item \textbf{Control y libertad del usuario.} La navegaci\'on, la salida de paneles y el reinicio de estado deben ser previsibles y reversibles.
    \item \textbf{Reconocimiento antes que memoria.} La interfaz debe reducir la necesidad de recordar nombres, ubicaciones o secuencias de interacci\'on, apoy\'andose en ayudas visuales y estados expl\'icitos.
\end{enumerate}

\subsubsection{Rotaci\'on mental y visualizaci\'on espacial}

Hegarty y Waller (2004) muestran que la interpretaci\'on de estructuras tridimensionales depende de habilidades espaciales que no se distribuyen homog\'eneamente entre usuarios. Bartlett y Dorribo Camba (2023) refuerzan que los modelos 3D interactivos pueden disminuir esta barrera al facilitar interpretaci\'on gr\'afica y comprensi\'on estructural. En consecuencia, el visor no se justifica solo por atractivo visual, sino como apoyo cognitivo para usuarios con distintos niveles de habilidad espacial.

\subsubsection{Percepci\'on visual y organizaci\'on de la informaci\'on}

Los principios de Gestalt (Wertheimer, 1923; K\"ohler, 1929) ayudan a estructurar la escena y la interfaz:
\begin{itemize}
    \item \textbf{Proximidad.} Elementos cercanos se perciben como pertenecientes al mismo subsistema.
    \item \textbf{Similitud.} El color o tratamiento visual de \textit{hotspots} y categor\'ias facilita asociaci\'on funcional.
    \item \textbf{Figura-fondo.} El modelo principal debe destacar sobre el fondo para no competir con el contenido t\'ecnico.
\end{itemize}

Ware (2012) y Miller (1956) tambi\'en sustentan el uso de categor\'ias crom\'aticas distinguibles, evitando esquemas que dificulten accesibilidad o saturen la memoria inmediata del usuario.

\subsubsection{Fundamentos matem\'aticos de optimizaci\'on gr\'afica}

\paragraph{Presupuesto poligonal.}
Una primera restricci\'on de complejidad geom\'etrica puede expresarse como:
\[
P_{\text{total}} = \sum_{i=1}^{n} P_i
\]
donde $P_i$ representa los tri\'angulos aportados por cada componente del modelo. En t\'erminos de planeaci\'on, el proyecto fija un presupuesto objetivo para mantener la escena dentro de un rango compatible con WebGL. Este presupuesto no es una ley universal del motor, sino un criterio de control para el proceso de optimizaci\'on.

\paragraph{Densidad de texel.}
La densidad de texel busca que la nitidez percibida de las texturas sea coherente entre componentes. Una forma operativa de expresarla es:
\[
\rho = \frac{R_{\text{texture}}}{L_{\text{mundo}}}
\]
donde $R_{\text{texture}}$ es la resoluci\'on efectiva aplicada sobre una direcci\'on o tramo de referencia, y $L_{\text{mundo}}$ es la longitud real correspondiente en el modelo. Si $\rho$ cambia bruscamente entre piezas, algunas se perciben borrosas y otras innecesariamente costosas en memoria.

\paragraph{Costo de dibujo y agrupaci\'on de mallas.}
El costo de renderizado no depende solo del n\'umero de tri\'angulos, sino tambi\'en del n\'umero de lotes o \textit{draw calls}. Sin agrupaci\'on, una aproximaci\'on simple es:
\[
D_{\text{base}} \approx \sum_{k=1}^{m} n_k
\]
donde $n_k$ representa la cantidad de objetos que requieren una combinaci\'on distinta de malla, material o pasada. Si existen condiciones de \textit{batching} o \textit{instancing} compatibles, varias instancias pueden reducirse a menos lotes efectivos:
\[
D_{\text{efectivo}} \approx \sum_{k=1}^{m} b_k, \qquad 1 \leq b_k \leq n_k
\]
Por tanto, la reducci\'on de \textit{draw calls} no debe modelarse como una divisi\'on exacta universal, sino como una funci\'on de compatibilidad entre materiales, mallas, pasadas de render y estado del \textit{pipeline}.

\paragraph{\textit{Frame time} y objetivo de FPS.}
En una aplicaci\'on interactiva, el tiempo por cuadro determina la fluidez percibida. Una aproximaci\'on \'util para perfilado es:
\[
T_{\text{frame}} \approx \max(T_{\text{CPU}}, T_{\text{GPU}})
\]
con un posible costo adicional por sincronizaci\'on. Para sostener 30 FPS, se requiere aproximadamente:
\[
T_{\text{frame}} \leq 33.33 \text{ ms}
\]
Esto traduce el rendimiento a una restricci\'on directamente medible en navegador y en herramientas de perfilado.

\subsubsection{Modelo matem\'atico de renderizado f\'isicamente basado (PBR)}

El renderizado f\'isicamente basado parte de la ecuaci\'on de render:
\[
L_o(\omega_o) = \int_{\Omega} f_r(\omega_i,\omega_o)\, L_i(\omega_i)\, (\mathbf{n}\cdot\omega_i)\, d\omega_i
\]
donde $L_o$ es la radiancia saliente en la direcci\'on de observaci\'on $\omega_o$, $L_i$ es la radiancia incidente desde la direcci\'on $\omega_i$, $\mathbf{n}$ es la normal de la superficie y $f_r$ es la BRDF del material. Esta ecuaci\'on expresa que el color visible depende de cu\'anta luz llega a la superficie, c\'omo la superficie la redistribuye y qu\'e parte de esa redistribuci\'on coincide con la direcci\'on de la c\'amara.

\paragraph{BRDF especular de Cook-Torrance.}
La componente especular puede modelarse como:
\[
f_{\text{spec}} = \frac{D(\mathbf{n},\mathbf{h})\, F(\mathbf{v},\mathbf{h})\, G(\mathbf{n},\mathbf{l},\mathbf{v})}{4(\mathbf{n}\cdot\mathbf{l})(\mathbf{n}\cdot\mathbf{v})}
\]
donde $\mathbf{l}$ es la direcci\'on de la luz, $\mathbf{v}$ la direcci\'on de vista y $\mathbf{h}$ el vector mitad entre ambas. Cada t\'ermino cumple una funci\'on distinta:
\begin{itemize}
    \item $D$ describe la distribuci\'on estad\'istica de microfacetas orientadas para reflejar luz hacia la c\'amara;
    \item $F$ modela el efecto Fresnel, es decir, el incremento de reflectancia cuando la observaci\'on se acerca al \'angulo rasante;
    \item $G$ aten\'ua la respuesta especular por auto-sombreado y enmascaramiento geom\'etrico entre microfacetas.
\end{itemize}

\paragraph{Aproximaci\'on de Fresnel de Schlick.}
Para reducir costo computacional, el t\'ermino de Fresnel suele aproximarse mediante:
\[
F = F_0 + (1 - F_0)(1 - \cos\theta)^5
\]
donde $F_0$ es la reflectancia a incidencia normal y $\theta$ es el \'angulo entre la direcci\'on de vista y el vector mitad. Esta aproximaci\'on permite obtener un comportamiento visualmente cre\'ible con un costo bajo en tiempo real.

\paragraph{Distribuci\'on GGX.}
La distribuci\'on de microfacetas adoptada com\'unmente en PBR en tiempo real es GGX:
\[
D_{\text{GGX}} = \frac{\alpha^2}{\pi\left((\mathbf{n}\cdot\mathbf{h})^2(\alpha^2 - 1) + 1\right)^2}
\]
donde $\alpha$ controla la amplitud del l\'obulo especular y suele derivarse de la rugosidad como $\alpha = roughness^2$. Valores bajos producen reflejos concentrados y valores altos reflejos anchos y difusos.

\paragraph{Separaci\'on difusa-especular y conservaci\'on de energ\'ia.}
En un flujo \textit{metallic-roughness}, la energ\'ia reflejada debe mantenerse acotada. De manera simplificada:
\[
k_d + k_s \leq 1
\]
donde $k_d$ representa la fracci\'on difusa y $k_s$ la especular. En materiales met\'alicos, la componente difusa tiende a desaparecer y el color especular toma protagonismo; en diel\'ectricos ocurre lo contrario.

\paragraph{Relaci\'on con la app real.}
En la aplicaci\'on final, el modo \texttt{Realistic} no reimplementa manualmente toda la BRDF microfacetada desde cero. El proyecto delega el c\'alculo PBR base al \textit{pipeline} de Unity URP mediante la funci\'on \texttt{UniversalFragmentPBR}, mientras que los shaders propios a\~naden l\'ogica espec\'ifica del sistema, como planos de corte globales y modos de inspecci\'on. Por ello, las ecuaciones anteriores se presentan como fundamento matem\'atico del modelo de iluminaci\'on empleado por URP y no como afirmaci\'on de un solver BRDF completamente escrito desde cero dentro del proyecto.

\paragraph{Iluminaci\'on basada en im\'agenes (IBL).}
La iluminaci\'on ambiental en PBR se apoya en mapas de entorno precalculados. Conceptualmente, la integral ambiental se aproxima mediante dos componentes: un mapa de irradiancia para la parte difusa y un mapa de entorno prefiltrado para la parte especular. Esto permite aproximar reflejos y contribuci\'on ambiental con costo constante por p\'ixel, sin necesidad de \textit{ray tracing} completo en tiempo real.

\subsubsection{Matem\'atica aplicada a shaders de inspecci\'on y an\'alisis}

\paragraph{Plano de corte global.}
El corte transversal de la app puede expresarse con la ecuaci\'on de un plano:
\[
\Pi = (n_x, n_y, n_z, d), \qquad d = -\mathbf{n}\cdot\mathbf{p}
\]
donde $\mathbf{n}$ es la normal del plano y $\mathbf{p}$ un punto perteneciente a dicho plano. Para cualquier punto del espacio $\mathbf{x}$, la distancia algebraica al plano se eval\'ua como:
\[
\mathbf{n}\cdot\mathbf{x} + d
\]
Si este valor es negativo, el fragmento queda del lado descartado del plano y no se renderiza:
\[
\mathbf{n}\cdot\mathbf{x} + d < 0 \Rightarrow \text{fragmento descartado}
\]
Esta formulaci\'on coincide con el mecanismo usado por los shaders del sistema cuando reciben los planos globales calculados por el gestor de corte.

\paragraph{Fresnel de borde para modos anal\'iticos.}
Modos como \textit{X-Ray} o \textit{Blueprint} usan una intensidad de borde basada en la relaci\'on entre normal y vista:
\[
F_{\text{borde}} = \left(1 - saturate(\mathbf{n}\cdot\mathbf{v})\right)^p
\]
donde $p$ controla la dureza del borde. Cuando la normal es casi perpendicular a la vista, la intensidad aumenta y se resaltan siluetas, contornos o regiones parcialmente transparentes.

\paragraph{Mapeo visual del modo t\'ermico.}
La vista t\'ermica del prototipo no corresponde a un solver FEA completo, sino a una visualizaci\'on h\'ibrida inspirada en un estado t\'ermico simplificado. A nivel visual, la temperatura normalizada puede expresarse como:
\[
T_{\text{vis}} = saturate\left(T_{\text{base}}\, s(\mathbf{x}) - e(\mathbf{x},\mathbf{v}) + \eta(\mathbf{x},t) + \delta_{\text{textura}}\right)
\]
donde $T_{\text{base}}$ proviene del subsistema t\'ermico, $s(\mathbf{x})$ representa el factor espacial radial o axial, $e(\mathbf{x},\mathbf{v})$ modela enfriamiento de borde, $\eta(\mathbf{x},t)$ introduce variaci\'on procedural suave y $\delta_{\text{textura}}$ incorpora microvariaciones visuales de la textura base. Esta ecuaci\'on describe una capa de visualizaci\'on, no una simulaci\'on termo-mec\'anica completa.

\newpage
